\chapter{Diarrhoea}
\index{Diarrhoea}

\section*{Definition and Overview}
Diarrhoea is the passage of loose or watery stools more often than usual. It is one of the most common illnesses worldwide and is usually caused by infection of the gut, but it can also result from food, medicines, and other conditions. Most episodes of diarrhoea are short-lived and resolve on their own within a few days, and the main treatment is replacing lost fluids. However, diarrhoea can be dangerous, especially in babies, young children, older adults, and people with weakened immunity, because it causes dehydration. This chapter covers the causes, symptoms, and management of diarrhoea.

\section*{Epidemiology}
Infectious diarrhoea affects virtually everyone at some point, with adults averaging around one episode a year and young children more. Gastroenteritis is one of the most common reasons for GP visits and for children missing school. Worldwide, diarrhoea is a leading cause of death in children under five, mainly in developing countries, although deaths have fallen with oral rehydration therapy and vaccines. In many countries, outbreaks occur in childcare, aged care, and from contaminated food, and travellers may develop diarrhoea abroad.

\section*{Aetiology and Pathophysiology}
\begin{itemize}
    \item \textbf{Viral infections:} rotavirus, norovirus, and adenovirus are the most common causes
    \item \textbf{Bacterial infections:} salmonella, campylobacter, shigella, and E. coli from contaminated food or water
    \item \textbf{Parasites:} giardia and cryptosporidium, often from contaminated water
    \item \textbf{Food poisoning:} toxins in contaminated food can cause rapid-onset diarrhoea and vomiting
    \item \textbf{Medicines:} antibiotics, laxatives, and some other drugs can cause diarrhoea
    \item \textbf{Other causes:} irritable bowel syndrome, inflammatory bowel disease, food intolerances, and malabsorption
    \item \textbf{Mechanism:} infection or irritation reduces the gut's ability to absorb fluid, so water is lost in the stool
\end{itemize}

\section*{Clinical Features}
\begin{itemize}
    \item Loose, watery, or frequent stools
    \item Abdominal cramps, bloating, and urgency
    \item Nausea, vomiting, and fever in infectious causes
    \item Dehydration: thirst, dry mouth, dark urine, reduced urine output, dizziness, and fatigue
    \item In severe dehydration: confusion, sunken eyes, rapid heart rate, and reduced skin elasticity
    \item Blood or mucus in the stool in some bacterial infections
    \item Most episodes settle within 2--5 days
\end{itemize}

\section*{Differential Diagnosis}
\begin{itemize}
    \item Viral versus bacterial gastroenteritis, which may need different management
    \item Food poisoning versus infection acquired from another person
    \item Irritable bowel syndrome, with chronic and recurrent loose stools
    \item Inflammatory bowel disease, with blood, mucus, and weight loss
    \item Coeliac disease and food intolerances, including lactose intolerance
    \item Medication-related diarrhoea
    \item Appendicitis or other abdominal emergencies, which can cause loose stools with pain
\end{itemize}

\section*{When to Seek Medical Care}
See a doctor if diarrhoea lasts more than a few days, is severe, or is accompanied by high fever, blood in the stool, or signs of dehydration. Babies, young children, older adults, and people with chronic illness need review sooner.

\textbf{Contact your local emergency services for:}
\begin{itemize}
    \item Signs of severe dehydration: confusion, drowsiness, fainting, or no urine for many hours
    \item Blood in the stool with severe abdominal pain or high fever
    \item A baby who is listless, has a sunken fontanelle, or is not passing wet nappies
    \item Vomiting that prevents keeping down any fluids
\end{itemize}

\section*{Diagnosis and Investigations}
\begin{itemize}
    \item \textbf{History and examination:} most cases are diagnosed clinically without tests
    \item \textbf{Stool tests:} for bacteria, parasites, or their toxins when infection is severe, prolonged, or in an outbreak
    \item \textbf{Blood tests:} electrolytes and kidney function in people with significant dehydration
    \item \textbf{Assessment of dehydration:} clinical signs guide treatment
    \item \textbf{Investigation of chronic diarrhoea:} tests for coeliac disease, inflammatory bowel disease, or intolerances
\end{itemize}

\section*{Management}
\begin{itemize}
    \item \textbf{Fluid replacement:} the most important treatment; oral rehydration solution is best for replacing fluid and electrolytes
    \item \textbf{Continue eating:} bland foods such as bread, rice, and bananas as tolerated
    \item \textbf{Avoid certain foods:} rich, fatty, and spicy foods may worsen symptoms
    \item \textbf{Medicines:} antidiarrhoeal agents are generally not recommended in infectious diarrhoea; they are used in some non-infectious cases
    \item \textbf{Probiotics:} may shorten some episodes of gastroenteritis
    \item \textbf{Antibiotics:} only for specific bacterial infections when indicated
    \item \textbf{Hospital care:} intravenous fluids for severe dehydration
    \item \textbf{Good hygiene:} hand washing and hygiene to prevent spread
\end{itemize}

\section*{Complications}
\begin{itemize}
    \item Dehydration and electrolyte disturbance, the most serious risks
    \item Kidney injury from fluid loss
    \item Malnutrition, especially in children with repeated episodes
    \item Spread of infection to family and community
    \item Prolonged symptoms or post-infectious irritable bowel syndrome
    \item Rare complications of specific infections, including haemolytic uraemic syndrome from E. coli
\end{itemize}

\section*{Prognosis and Outcomes}
Most episodes of diarrhoea resolve within a few days without treatment other than fluids. With appropriate rehydration, complications are uncommon in otherwise healthy people, and death from diarrhoea is rare in many countries. Children and older adults recover well when dehydration is prevented and treated early. Chronic causes of diarrhoea respond to treatment of the underlying condition.

\section*{Prevention and Self-Care}
\begin{itemize}
    \item Wash hands thoroughly with soap and water, especially after the toilet and before food
    \item Practise safe food handling, storage, and cooking
    \item Drink safe water, including when travelling
    \item Keep children up to date with rotavirus vaccination
    \item During illness, drink plenty of fluids and use oral rehydration solution
    \item Stay home while infectious to avoid spreading the illness
    \item Seek early medical advice for babies, older adults, and people with chronic illness
\end{itemize}

\section*{Sources and Further Reading}
The following reputable sources provide further information for healthcare students and patients seeking reliable, current advice about diarrhoea.

\begin{itemize}
    \item Healthdirect Australia. \textit{Diarrhoea: causes and treatment.} https://www.healthdirect.gov.au/
    \item NSW Health. \textit{Gastroenteritis and diarrhoea fact sheets.} https://www.health.nsw.gov.au/
    \item NHS. \textit{Diarrhoea: causes, treatment, and when to see a doctor.} https://www.nhs.uk/conditions/diarrhoea/
    \item Mayo Clinic. \textit{Diarrhoea: symptoms and treatment.} https://www.mayoclinic.org/
    \item MSD Manual. \textit{Diarrhoea: professional overview.} https://www.msdmanuals.com/
    \item World Health Organization. \textit{Diarrhoeal disease fact sheet.} https://www.who.int/
\end{itemize}
